%% ma: L12-B14-0064
%% lop: 12
%% chuong: 5
%% bai: 14
%% dang: 12.14.21
%% loai: TN4
%% muc: TH
%% dap_an: "C"
%% nguon: "To_n_12_PT_M_t_Ph_ng.pdf (Đề kiểm tra đánh giá kiến thức, Chuyên đề Phương trình mặt phẳng, Nguyễn Anh Khôi - Lê Minh Khôi), Câu 19"
%% kien_thuc: "Bài 14 (phương trình mặt phẳng theo đoạn chắn)"
%% trang_thai: cho-duyet
%% ly_do: "Dạng 12.14.21 mới đề xuất, chờ giáo viên duyệt."
%% ghi_chu: "Tài liệu gốc không kèm đáp án; đáp án do Claude giải và kiểm bằng sympy. Đề gốc cho phương trình $ax+by+cz+d=0$ và hỏi $a+b+c+d$: không xác định vì nhân cả phương trình với $k$ thì tổng thành $k$ (cả A và C đều có thể đúng). Đã sửa đề, cố định $d=-12$; đáp án C như ý đề gốc."
\begin{ex}
Trong không gian $Oxyz$, cho ba điểm $A(2;0;0)$, $B(0;3;0)$, $C(0;0;4)$. Biết phương trình mặt phẳng $(ABC)$ là $ax+by+cz-12=0$. Khi đó giá trị của $a+b+c-12$ bằng
\choice{$-1$}{$9$}{\True $1$}{$-9$}
\loigiai{$(ABC)\colon \dfrac{x}{2}+\dfrac{y}{3}+\dfrac{z}{4}=1\Leftrightarrow 6x+4y+3z-12=0$. Suy ra $a=6$, $b=4$, $c=3$ và $a+b+c-12=1$. Chọn C.}
\end{ex}
